% Section 5.1 -- vocabulary coverage. The headline result.
% Ported from the frozen MATH-AI results section and expanded: intersection
% counts alongside every Jaccard so each row is recomputable, and the fifth
% design (common support) that four pages had no room for.
% No literal numbers: every figure is a macro from numbers.tex.
% The uncorrected human count is retracted and has no macro by design.

\subsection{Vocabulary coverage}
\label{sec:coverage}

Over the full corpora the two vocabularies intersect in \intersectionFull{}
premises, giving a Jaccard index of \jaccardFull{}. \humanOnlyFullPct{}\% of
the premises mathlib's own proofs cite are never cited by any machine proof in
the corpus. The asymmetry
is close to containment. Only \machineOnlyFull{} premises are machine-only, so
the machine vocabulary is very nearly a subset of the human one rather than a
different region of the library. Two properties of that small vocabulary are
reported in the appendix. It is weighted toward language infrastructure rather
than mathematics, at \machineVocabInfraPct{}\% against the human corpus's
\humanVocabInfraPct{}\%, leaving a mathematical working vocabulary of
\machineVocabMath{} declarations (Appendix~\ref{sec:layer-composition}). And
it is used very unevenly, with proofs closing in at most two steps under an
arithmetic tactic set making up \degenerateMachineSharePct{}\% of the machine
corpus against \degenerateHumanSharePct{}\% of mathlib
(Appendix~\ref{sec:concentration}). Every uncorrected instrument asymmetry we
have identified narrows the gap reported here rather than widening it, so these
figures are floors; Appendix~\ref{app:limitations}, L2, enumerates them.

The designs differ in what they are allowed to explain away. The full corpora
control for nothing. Rarefaction removes the corpus-size advantage by drawing
the human corpus down to the machine corpus's proof count. Common support
discards the cells of the (tactic class $\times$ tactic step) grid that only
one corpus occupies, so no rate is imputed where the other side has no
population. The area-restricted design keeps only the regions of the library
the machine vocabulary occupies at all, removing the objection that the two
corpora are simply doing different mathematics. The matched design is the
strictest: equal proof counts, equal tactic class, equal step bin.

\begin{table}[t]
\centering
\small
\begin{tabular}{lrrrrrr}
\toprule
design & human & machine & ratio & Jaccard & $|H \cap M|$ & human-only \\
\midrule
full corpora & \humanVocab{} & \machineVocab{} & \vocabRatio{}$\times$ & \jaccardFull{} & \intersectionFull{} & \humanOnlyFullPct{}\% \\
\textbf{area-restricted} & \textbf{\humanVocabAreaRestricted{}} & \textbf{\machineVocab{}} & $\mathbf{\areaRestrictedRatio\times}$ & \textbf{\areaRestrictedJaccard{}} & \textbf{\intersectionAreaRestricted{}} & \textbf{\areaRestrictedHumanOnlyPct{}\%} \\
1:1 class- and step-matched & \humanVocabMatched{} & \machineVocabMatched{} & \vocabRatioMatched{}$\times$ & \jaccardMatched{} & \intersectionMatched{} & \humanOnlyMatchedPct{}\% \\
\bottomrule
\end{tabular}
\caption{Premise vocabulary under three designs of increasing severity. Each
row carries its intersection count, so its Jaccard index is recomputable from
the row alone as $|H \cap M| / (|H| + |M| - |H \cap M|)$. The area-restricted
row is the primary control (Section~\ref{sec:size-controls}); the matched row
is the strictest design we ran (Section~\ref{sec:size-controls}). Two further
designs are reported in the text rather than here: rarefying the human corpus
to the machine corpus's size gives \vocabRatioRarefied{}$\times$
(Section~\ref{sec:size-controls}), and restricting to common support gives
\vocabRatioSupport{}$\times$ at a Jaccard of \jaccardSupport{}, against
\jaccardFull{} on the full corpora.}
\label{tab:coverage}
\end{table}



The point of the table is the invariance rather than any single row. The ratio
never leaves the range \vocabRatioMatched{}$\times$ to \vocabRatio{}$\times$,
and the human-only fraction never falls below \humanOnlyMatchedPct{}\%, under
controls that variously remove corpus size, proof length, tactic style and
subject area. The magnitude is design-dependent and we quote the range rather
than the largest figure; the direction and order of magnitude are not.
Coverage is also the most stable quantity we measure. Moving from the full
corpora to common support changes the Jaccard index only in its fourth decimal
place, \jaccardFull{} to \jaccardSupport{}, while that restriction discards
roughly a quarter of the machine corpus. The result is therefore not an
artifact of which proofs happen to be comparable. Every proof on both sides of
this comparison is verified. Correctness is satisfied throughout, and none of
the structure in this table is visible to \texttt{pass@k}.
